% Regression: t2_sptintertwin — imported current-behavior fixture from the handoff corpus.
% Formula: (per site, MPO virtual spaces):
% T2 — RMP arXiv:2011.12127, fig. fig3-pepe_sptintertwin-gh: the SPT MPO
% intertwiner (beyond hard09_v2: X_gh now renders as the source's tall
% labeled BOX, not a fusion bar).
% Formula (per site, MPO virtual spaces):
%     u_gh(1) . X_gh . [u_g (x) u_h](2) [u_g (x) u_h](3)
%   = u_gh(1) u_gh(2) . X_gh . [u_g (x) u_h](3),
% i.e. X_gh : V_gh -> V_g (x) V_h intertwines the gh MPO with the stacked
% g,h MPOs.  Reference ink: gh / g / h are MPO tensors (red virtual wire
% through the box, black physical wire through it vertically); X_gh is a
% tall white box with ONE red wire west and TWO red wires east.
%
% Declarative encoding: X_gh = \tn[box, wires=2, no legs] in a two-row
% operator picture -- the wires=k atom terminates both east wires ON the
% box and overshoots the rows exactly like the source's tall box; the
% picture-level `west=none` seals the box's west face, and the single
% combined west wire is supplied by the juxtaposed one-row gh picture
% (its open east stub abuts the box).  The g/h columns share one vertical
% physical wire through both rows (op:operator pairing).
% Noted gap (unchanged since hard09): no cell can sit ON a combined wire
% inside one picture, so the gh factor remains a separate juxtaposed
% picture and the contraction at the seam is implied, not drawn; the
% two-picture split also splits the boundary signature.
\documentclass[varwidth=19cm, border=6pt]{standalone}
\usepackage{amsmath, amssymb}
\usepackage{tenkz}
\begin{document}
\[
  % LHS: gh at site 1, then X_gh, then (g over h) at sites 2 and 3
  \begin{tenkz}[rows={op:operator}, tensor style=box]
    \tn{gh}
  \end{tenkz}
  \!
  \begin{tenkz}[rows={op:operator, op:operator}, west=none,
                tensor style=box]
    \tn[box, wires=2, no legs]{X_{gh}} & \tn{g} & \tn{g}\\
                                       & \tn{h} & \tn{h}
  \end{tenkz}
  \;=\;
  % RHS: gh at sites 1 and 2, then X_gh, then (g over h) at site 3
  \begin{tenkz}[rows={op:operator}, tensor style=box]
    \tn{gh} & \tn{gh}
  \end{tenkz}
  \!
  \begin{tenkz}[rows={op:operator, op:operator}, west=none,
                tensor style=box]
    \tn[box, wires=2, no legs]{X_{gh}} & \tn{g}\\
                                       & \tn{h}
  \end{tenkz}
\]
\end{document}
